\documentclass[11pt]{article}
\usepackage[T1]{fontenc}
\usepackage{lmodern,microtype,amsmath,amssymb,amsthm,booktabs}
\usepackage[margin=1in]{geometry}
\usepackage[hidelinks]{hyperref}
\hypersetup{pdftitle={Compatible local replacements for shifted lonely runners},
           pdfauthor={Chengyi Fang}}
\setlength{\emergencystretch}{1em}
\allowdisplaybreaks
\newtheorem{theorem}{Theorem}
\newtheorem{lemma}[theorem]{Lemma}
\newtheorem{proposition}[theorem]{Proposition}
\theoremstyle{definition}
\newtheorem{definition}[theorem]{Definition}
\newcommand{\T}{\mathbb{R}/\mathbb{Z}}
\newcommand{\Z}{\mathbb{Z}}
\newcommand{\norm}[1]{\left\lVert #1\right\rVert}
\newcommand{\fl}[1]{\left\lfloor #1\right\rfloor}
\DeclareMathOperator{\lcm}{lcm}
\title{Compatible local replacements for shifted lonely runners}
\author{Chengyi Fang\\[-2pt]\small
\href{https://orcid.org/0009-0006-1237-7763}{ORCID 0009-0006-1237-7763}}
\date{26 November 2026}

\begin{document}
\maketitle
\begin{abstract}
Write \(1/(n+1+E_n)\) for the infimal covering radius of \(n\) shifted
runners with distinct positive integer speeds. Combining three published
local rules, with shared runners where the prescribed phases agree, gives
\[
 E_n\ge \max\left\{\fl{n/150},
          \fl{\frac{8011n-3241770}{1085489}}\right\}\qquad(n\ge1).
\]
An additional family of disjoint replacements yields
\(\liminf_{n\to\infty}E_n/n\ge 513079/69470921\).
We also give a finite decision procedure for affine local rules with
unrestricted odd added speeds and real phase parameters.
\end{abstract}

\section{Covering radius and local replacements}

Let \(\norm{x}=\min_{k\in\Z}|x-k|\). For distinct positive integers
\(v_1,\ldots,v_n\) and phases \(s_i\in\T\), put
\[
 \gamma(v,s)=\max_{t\in\T}\min_{1\le i\le n}\norm{v_it+s_i},
 \qquad \Gamma(n)=\inf_{v,s}\gamma(v,s),
 \qquad E_n=\frac1{\Gamma(n)}-(n+1).
\]
A radius \(d\) covers the time circle when the intervals on which some
runner has distance at most \(d\) cover \(\T\). The stationary observer
is excluded from \(n\). A single runner covers a fraction \(2d\) of the
circle for \(0\le d\le 1/2\), so \(\Gamma(n)\ge 1/(2n)\).
Blanco, Criado and Santos proved \(\Gamma(n)<1/(n+1)\) for
\(5\le n\le17\) \cite[Theorem 1.8]{BCS}. Poliakova constructed
counterexamples for every \(n\ge95\) and proved \(E_n\ge\fl{n/287}\)
\cite[Theorems 1 and 2]{Poliakova}. Her strongest coefficient in
Table 1 is \(1/224\). The construction below improves the quantitative
bound by combining her rules across scales and accounting for overlaps.
Zhang's question about every \(n\ge5\) and the global minimizers
\cite[Question 8.5]{Zhang} also motivates exact computations at
individual sizes \cite{Upstream}.

Set \(\delta=1/(N+1)\). The zero-phase speeds \(2,4,\ldots,2N\) give
a \(\delta\)-cover. Of the \(N+1\) points \(0,2t,\ldots,2Nt\) modulo
one, two coincide or two successive points in circular order have
separation at most \(\delta\). Their index difference supplies a runner
covering \(t\). Call the smallest covering speed responsible for \(t\),
following Poliakova \cite[Definition 5]{Poliakova}.

A rule consists of a set \(R\) of deleted even speeds and a table \(P\)
of speeds \(v\) with prescribed phases \(a_v+h_v\delta\). The even entries
of \(P\) are rephased seed speeds and the odd entries are additions.
An even entry may be prescribed the identically zero phase.
The table has distinct speeds and \(R\cap P=\varnothing\), where a speed
in \(P\) means a speed in its domain. The affected even set is
\(Q=R\cup P_{\mathrm{even}}\), the support is \(Q\cup P\), and the gain
is \(|R|-|P_{\mathrm{odd}}|\).
At an odd scale \(b\), multiply each speed by \(b\) and retain its
phase formula. Such an application fits as soon as \(N\ge b\max Q/2\).

\begin{lemma}[Primitive-residue test]\label{lem:primitive}
For the rule to cover every time for which a speed it affects was
responsible in the even seed, for every \(N\) and admissible odd scale,
it is necessary and sufficient that the intervals
\[
 \left[\frac qv(-1-h_v),\,\frac qv(1-h_v)\right],
 \qquad v\in P,\quad vk/q+a_v\in\Z ,
 \tag{1}\label{eq:test}
\]
cover \([-1,1]\) for every \(q\in Q\) and \(0\le k<q\)
with \(\gcd(k,q/2)=1\).
\end{lemma}
\begin{proof}
If \(qb\) is responsible at \(t\), write \(qbt=k+\delta z\) with
\(k\in\Z\) and \(|z|\le1\). An integer divisor \(r>1\) common to \(k\)
and \(q/2\) would give a smaller even seed speed \(qb/r\) with distance
at most \(\delta/r\). Thus the residue of \(k\) modulo \(q\) is among
those tested. For each admissible \(v\) in \eqref{eq:test},
\(vbt+a_v+h_v\delta\) differs from an integer by
\(\delta((v/q)z+h_v)\). This proves sufficiency.

For necessity take \(b=1\) and fix any tested \(q,k\) and any
\(z\in[-1,1]\). At \(k/q\) each smaller even speed has positive distance
from the integers by the coprimality condition. At
\(t=(k+\delta z)/q\) they still have distance greater than \(\delta\)
for all sufficiently large \(N\), whereas \(q\) covers. Each replacement
with \(vk/q+a_v\notin\Z\) also has distance greater than \(\delta\)
there for sufficiently large \(N\).
For an aligned replacement and sufficiently small \(\delta\), the
distance is exactly \(\delta|(v/q)z+h_v|\).
Locality forces one aligned replacement's interval to contain \(z\).
\end{proof}

A collection of rules with disjoint supports composes. At each time,
an unaffected responsible speed still covers, and an affected one has a
witness of its own rule in the combined configuration.
There is also an overlapping composition. Take the union of all prescribed
speeds when each shared speed has the same affine phase (the constant term
is read modulo one). Remove only those speeds listed for deletion that
are absent from the prescriptions. Every affected speed was affected by
some member rule, and all witnesses of that rule survive. Hence its local
coverage survives. We count the gain after the union and the removal
of duplicates. A half-turn translation in time preserves the zero
phases of the even seed and adds \(1/2\) to the phase of each odd speed.
It preserves locality for the translated rule, since responsibility for
the seed is unchanged. This translation acts in the same way after
multiplication by any odd scale.

\section{Three rules and their combinations}

Use rule \(A\) from Poliakova's Theorem 7 and rules \(B,C\) from the
\(19\to17\) and \(24\to21\) rows of her Table 1 \cite{Poliakova}.
For the exact data below, \([X]_{a,h}\) prescribes phase
\(a+h\delta\) at each speed in \(X\). All additions are odd.
\begin{align*}
 R_A={}&\{8,32,192\},\\
 P_A={}&[\{12,24,48,96\}]_{1/2,\,1/2}
       \cup[\{16,64\}]_{0,\,-1}\cup[\{3,9\}]_{1/4,\,0},\\[2pt]
 R_B={}&\{8,32,36,40,72,180,192,240\},\\
 P_B={}&[\{12,24,48,60,96,108,120\}]_{1/2,\,1/2}
       \cup[\{16,64,80\}]_{0,\,-1}\cup[\{216\}]_{0,\,-2}\\
     &{}\cup[\{3,9,27\}]_{1/4,\,1/4}
       \cup[\{15,45,135\}]_{3/4,\,1/4},\\[2pt]
 R_C={}&\{4,16,20,36,40,64,72,240,384,432\},\\
 P_C={}&[\{12,24,48,96,192\}]_{1/2,\,1/2}
       \cup[\{60,108,120,216\}]_{1/2,\,-1/2}\\
     &{}\cup[\{8,128\}]_{0,\,-1}\cup[\{80,144\}]_{0,\,1}
        \cup[\{32\}]_{0,\,-3/2}\\
     &{}\cup[\{3,45,81\}]_{1/4,\,-1/4}
        \cup[\{5,9\}]_{1/4,\,1/4}
        \cup[\{15,27\}]_{1/4,\,0}.
\end{align*}
For example the shared prescriptions of \(B_1\) and \(B_5\) are at
\(15,45,60,80,120,135\). Time-translating the second copy by \(1/2\)
makes all six agree, and speed \(240\) is retained with its prescribed
phase. This merged rule deletes thirteen even speeds and adds nine
odd speeds. The merge \(A_1,C_5,C_9\) has shared prescriptions at
\(45,135,405,540,720,1080\), all with the same phase in their respective
copies. It deletes twenty even speeds and adds thirteen odd speeds.

A subscript denotes the odd scale and \(B_5'\) denotes the
time-translated copy. The following counts come from applying (1) to
the combined tables. The threshold is half the largest affected even
speed. The final column counts the pairs \((q,k)\) checked before
any further scaling.
\begin{center}
\begin{tabular}{@{}rlrr@{}}
\toprule
Threshold & Rule combination & Gain & Residues \\
\midrule
96 & \(A_1\) & 1 & 184 \\
120 & \(B_1\) & 2 & 536 \\
216 & \(C_1\) & 3 & 882 \\
600 & \(B_1,B_5'\) & 4 & 2\,680 \\
1080 & \(A_1,A_9,C_5\) & 5 & 5\,356 \\
1944 & \(A_1,C_5,C_9\) & 7 & 9\,928 \\
5832 & \(A_1,A_5,C_{25},C_{27}\) & 8 & 41\,576 \\
9720 & \(C_1,C_{25},A_{45},B_{81}\) & 9 & 65\,766 \\
\bottomrule
\end{tabular}
\end{center}
To reproduce a check, merge the prescriptions and remove the
prescribed speeds from the deletion set. For each \(q\) and \(k\)
from Lemma~\ref{lem:primitive}, form the eligible intervals and sort by
left endpoint. Start at \(-1\), and advance the rightmost covered point
to the greater of its present value and the next right endpoint if the
next interval starts at or below it. The check succeeds when the
rightmost point reaches \(1\). This uses rational endpoints and accepts
coincidences. The full standard-library check is
\texttt{overlap\_packing.py}.

All entries of these supports are \(2,3,5\)-smooth. If \(c\) is coprime
to \(30\), scaling a combination by \(c\) gives a rule whose speeds have
prime-to-\(30\) part \(c\). Here this part is the integer left after
dividing away every factor \(2,3,5\). Thus supports at different cores
cannot meet. Within one core we choose a single combined rule and count
each prescribed speed once. Every remaining even seed speed then occurs
once and every added speed is odd. These facts ensure positivity and
distinctness throughout the construction.

\section{Counting the savings}\label{sec:count}

For \(x\ge0\), let \(g(x)\) count reached thresholds in the multiset
\[
  (96,120,216,600,1080,1944)
     \ \sqcup\ \bigsqcup_{j\ge0}5^j(1944,5832,9720).
 \tag{2}\label{eq:thresholds}
\]
At a core \(c\) use \(x=N/c\). For \(x<9720\) the table realizes
\(g(x)\), including the jump of two at \(1944\).
The supports of \(A_{9^i}\) are pairwise disjoint for all \(i\ge0\).
Indeed, in the support of \(A\), the only multiple of \(9\) is \(9\)
and \(1\) is absent. These supports are also prime to \(5\).
For \(x\ge9720\) all six head thresholds and all three first-cycle
thresholds have been reached. All six head thresholds have been reached
at \(x/5\). Consequently \(g(x)=3+g(x/5)\).
Use the recipe at \(x/5\) scaled by \(5\) together with the three disjoint
rules \(A_1,A_9,A_{81}\), which fit at \(x\ge9720\) because their highest
threshold is \(96\cdot81=7776\). This realizes \(g(x)\) by induction.
At \(x=9720\) the recursion attains gain ten, accounting for two reached
thresholds with this same value.

Write \(C_m(x)=\#\{c\in\mathbb Z_{>0}:c\le x,\gcd(c,m)=1\}\).
The construction saves \(G_0(N)=\sum_{\gcd(c,30)=1}g(N/c)\) runners.
Each integer coprime to \(6\) equals \(5^j c\) uniquely with
\(\gcd(c,30)=1\). Thus
\[
 G_0(N)=\sum_{t\in (96,120,216,600,1080,1944)}C_{30}(N/t)
        +\sum_{t\in (1944,5832,9720)}C_6(N/t).
 \tag{3}\label{eq:period}
\]
The two appearances of \(1944\) contribute separately.

For the larger family let \(g_1(x)=g(x)\) if \(x<9720\), and
otherwise use the improved recipe from \(x/5\) scaled by five with
every \(A_{9^i}\) satisfying \(96\cdot9^i\le x\).
The support argument above still applies. Subtracting the recursion
for \(g\) from that for \(g_1\) shows
\[
 g_1(x)-g(x)=\sum_{i\ge3,\ j\ge0}
   \mathbf{1}_{96\cdot9^i5^j\le x}.
\]
Each added threshold is above \(9720\) and only finitely many are reached
at a fixed \(x\). Therefore this construction saves
\[
 G_1(N)=G_0(N)+\sum_{i\ge3} C_6\bigl(N/(96\cdot9^i)\bigr).
 \tag{4}\label{eq:extended}
\]
In particular \(G_1(N)\ge G_0(N)\). Equations
\eqref{eq:period}--\eqref{eq:extended} give an exact count of the
savings of the specified constructions. They suffice to certify
larger excess for any particular \(n\) if the number saved at
\(N=n+h\) is at least \(h\). In that event the constructed cover has
at most \(n\) runners, and adding unused positive integer speeds at
phase zero if necessary gives \(\Gamma(n)\le1/(n+h+1)\).

\begin{theorem}\label{thm:main}
For every \(n\ge1\),
\[
 E_n\ \ge\ \max\!\left\{
        \fl{\frac n{150}},
        \fl{\frac{8011n-3241770}{1085489}}\right\}.
 \tag{5}\label{eq:main}
\]
Moreover
\[
  \liminf_{n\to\infty}\frac{E_n}{n}\ \ge\
       \frac{513079}{69470921}.
 \tag{6}\label{eq:liminf}
\]
\end{theorem}
\begin{proof}
For the first floor use \(h=\fl{n/150}\) and \(N=n+h\ge151h\).
If \(h=0\) use the even seed. Otherwise it suffices to show
\(G_0(151h)\ge h\). The residues of positive integers coprime to
\(30\) are \(1,7,11,13,17,19,23,29\) modulo \(30\). Counting
complete periods and checking the \(30\) remainders gives
\(C_{30}(x)\ge4x/15-14/15\) for \(x\ge0\).
Using the first seven thresholds in \eqref{eq:thresholds} gives
\[
  G_0(151h)\ \ge\ \frac{792599}{729000}h-\frac{98}{15}>h
  \qquad (h\ge75).
\]
For \(1\le h\le74\), the sum over the first nine thresholds minus \(h\)
has minimum zero. This finite check uses
\(C_{30}(x)=8\fl{x/30}+\#\{r\in\{1,7,11,13,17,19,23,29\}:
r\le\fl{x}\bmod30\}\).

For the second floor set \(p=8011,\ q=1093500\).
The function \(H(N)=pN-qG_0(N)\) on nonnegative integers has period
\(P=17496000\). Indeed \(P\) is the least common multiple of \(mt\) for
the nine terms \(C_m(N/t)\) in \eqref{eq:period}. Over a full period
the sum increases by \(128176=pP/q\). The largest value of \(H\) is
\(K=4335269\), attained at \(N=8139479\) where \(G_0(N)=59626\).
For this finite calculation initialize \(G_0(0)=0\) and enumerate
the progressions of jumps \(t(r+jm)\) with \(1\le r<m\), \(\gcd(r,m)=1\)
and \(j\ge0\), stopping at \(P\). Add the multiplicities of equal jumps.
Between jumps \(H\) is increasing, so check the integer immediately
before each of the \(125176\) distinct jump positions and both
period endpoints. This yields \(K\) by exact integer comparisons.
The accompanying verifier \texttt{periodic\_packing.py} also
evaluates the maximizing case directly by quotient and remainder.

Set \(h=\fl{(pn+q-1-K)/(q-p)}\). When \(h>0\),
\(p(n+h)-K\ge q(h-1)+1\), so \(G_0(n+h)>h-1\). Integrality yields
\(G_0(n+h)\ge h\), as required. The case \(h\le0\) follows from the
even seed.

Finally \(G_0(N)/N\to p/q\) by \eqref{eq:period}. For the additional
terms of \(G_1(N)/N\), the bound \(C_6(N/t)\le N/t\) gives a
summable majorant, and \(C_6(N/t)/N\to1/(3t)\) termwise.
Hence
\[
 \lim_{N\to\infty}\frac{G_1(N)}N
   =\frac{8011}{1093500}
       +\sum_{i\ge3}\frac1{3\cdot96\cdot9^i}
   =\frac{8011}{1093500}+\frac1{186624}
   =\frac{513079}{69984000}=\alpha_1.
\]
Given \(0<a<\alpha_1\), for all large \(n\) take
\(h=\fl{an/(1-a)}\) and \(N=n+h\). Then \(G_1(N)\ge aN\ge h\).
Letting \(a\) increase to \(\alpha_1\) proves
\eqref{eq:liminf}.
\end{proof}

\paragraph{Disjoint compositions.}
The first six disjoint-support thresholds are
\(96\), \(120\), \(216\), \(864\), \(1080\), and \(3000\).
They are followed by \(5^j(5400,5832,9720)\)
for \(j\ge0\). For the distinct intermediate gains \(4,5,6,7\)
one can use respectively
\[
 \begin{gathered}
 (A_1,B_5,A_9),\qquad (A_1,A_9,C_5),\\
 (A_1,A_5,B_9,B_{25}),\qquad (A_1,A_5,B_9,C_{25}).
 \end{gathered}
\]
The other gains through nine use the same disjoint rows as above.
The same three-rule extension gives the repeated tail.
The saved fraction is \(76897/10935000\), yielding
\(\liminf E_n/n\ge76897/10858103\). The disjoint-only bound is
\(E_n\ge\fl{n/167}\). Indeed use \(N=168h\) with \(h=\fl{n/167}\);
the first three contributions sum to at least
\((707/675)h-14/5>h-3\). When \(h\ge7\) the fourth and fifth each
contribute at least one. For \(h=1,\ldots,6\), the first nine
contributions sum to \(2,3,3,4,5,6\), respectively.
The disjoint and overlap savings are independent constructive bounds.

\section{A finite decision problem for affine local rules}\label{sec:finite}

Fix an even integer \(L\ge2\) and a target gain \(r\ge1\). Consider
exactly the class of Lemma~\ref{lem:primitive}, in which the affected
even speeds do not exceed \(L\), the added speeds are positive odd
integers, and each phase is affine in \(\delta\) with coefficients
independent of \(N\).

\begin{proposition}\label{prop:np}
The existence of such a rule of gain at least \(r\) is decidable.
With \(L\) in unary and \(r\) given as an integer, the decision
problem belongs to NP.
\end{proposition}
\begin{proof}
Fix the affected set \(Q\) of any viable rule and let \(M=\lcm(Q)\).
The number of odd additions is at most \(L/2-r\).
Within each residue class modulo \(M\) compress the added speeds,
in order, to the smallest distinct positive members of that class.
Since \(M\) is even, these speeds are odd. Replacing \(v\) by \(v'\le v\),
keep \(a_v\) and replace the slope by \(h_vv'/v\).
Alignment for every tested \(q,k\) is unchanged, while the absolute
value of the error expression in (1) is multiplied by \(v'/v\).
Thus all obligations remain covered and all added speeds are at
most \((L/2-r)M\).

Discard unused additions and restore unused even prescriptions to
phase zero. The latter can only remove affected rows. Each remaining
constant phase appears in some alignment, so it has denominator
dividing the original \(M\) when normalized modulo one.
Write \(u_v=-h_v/v\). On each tested row with affected speed \(q\)
the eligible intervals in the coordinate \(z/q\) have centers \(u_v\)
and radii \(1/v\), and must cover \([-1/q,1/q]\).
Choose a simple intersecting chain of eligible intervals whose union
covers the target. Such a chain exists by restricting the interval
graph to intervals meeting the target and connecting one that contains
each endpoint. Every chain constraint is a difference inequality on
centers or on one center and the reference \(u_0=0\), with bound a sum
of signed terms from \(1/q,1/v\).
We can discard or restore any prescriptions unused by these chains.

Let \(D\) be the lcm of the original \(Q\) together with the remaining
prescribed speeds. Scale all centers by \(D\). The right-hand sides
are integers, so flooring every coordinate preserves every selected
inequality, since \(x_i\le x_j+c\) implies
\(\lfloor x_i\rfloor\le\lfloor x_j\rfloor+c\) for integral \(c\).
The reference stays zero. Each originally used center satisfies
\(|u_v|\le1/2+1/v\) because its chain interval meets a target segment.
The resulting certificate has \(u_v\in D^{-1}\Z\), with
\(|u_v|\le1/2+1/v+1/D\).

One can encode the certificate using the universal \(M\) equal to
\(\lcm(2,4,\ldots,L)\) and the corresponding larger grid.
By a product bound its bit length is \(O(L\log L)\).
There are at most \(L/2\) prescribed speeds, so
\(\log D=O(L^2\log L)\). The checker tests the gain and
distinctness and verifies at most \(\sum_{q\le L}q\) interval unions
with polynomial bit complexity. No chain choices need be stored.
\end{proof}

\paragraph{A complete search without scanning a phase grid.}
Include every even \(q\le L\) in the rows, including retained zero-phase
speeds. A row is indexed by \(r_0=k/q\) from
Lemma~\ref{lem:primitive} and has target \([-1/q,1/q]\).
Identical fractions can be merged by keeping the largest target
interval. A retained even speed with \(a_v=h_v=0\) covers each of its
own rows. Conversely any cover of these rows by at most \(L/2-r\)
runners gives a rule as in Lemma~\ref{lem:primitive}, with each unused
even speed deleted. Declare all kept even speeds through \(L\) to be
prescribed, including the ones with zero phase. The assigned
witnesses all lie in the finite cover.

Start each row at its left endpoint. Suppose some rows are unfinished,
and pick a row with the smallest covered frontier \(f\). In any
completion some unused eligible runner covers an interval immediately
to the right of \(f\) on this row, because the existing covered
segments outside the frontier are empty in this starting and updated
representation. A finite completion has a member covering an entire
right neighborhood of \(f\), with left endpoint at most \(f\).
Place this runner's left endpoint at \(f\), preserving its width
\(2/v\) and its alignment on the chosen row. The position requires
\(a_v=-vr_0\bmod1\) and \(h_v=-1-vf\). It loses only points strictly
to the left of \(f\), already covered on every unfinished row where
the runner was eligible. Eligibility on another row \(s_0\) is
equivalent to \(\operatorname{den}(s_0-r_0)\mid v\).
On eligible rows we advance the frontier to \(f+2/v\) if this is an
increase, capped at the right endpoint. These advances maintain
connected coverage from the original left endpoint since the leftmost
frontier was selected. On every unfinished row, all previously placed
intervals met its original or then-current covered segment whenever
they extended it, so they leave no covered points beyond its current
frontier. Branching on the unused eligible speeds gives soundness and
completeness by induction on the remaining runner budget.

Only finitely many speeds need be branched on. Put
\(M=\lcm(2,4,\ldots,L)\). Two odd speeds with the same
\(\gcd(v,M)\) allow precisely the same row incidences after alignment
at a common row, since each relevant denominator divides \(M\).
Within any class choose at most the first \(L/2-r\) positive speeds.
To see completeness, replace the speeds of a completion by the
smallest distinct speeds in that class, sorted in matching order.
Align each replacement at an original served row and retain its
interval's left endpoint. The interval grows and all incidences are
preserved. Even speeds are left unchanged.

The exact search \texttt{local\_pivot.py} finds no cover of these rows
with \(L/2-1\) runners for \(L=2,4,6,8,10,12\). At \(L=12\) it
uses 24 rows and 26 candidate speeds. Its pruning is also exact.
At any step a set of pairwise incompatible active rows requires as
many remaining runners as there are rows in the set. When just
one runner remains, it must align with every active row and its
right endpoint must reach the largest active target endpoint.
This excludes positive gain through \(L=12\) in the stated class.
The cutoff is a limit on local even speeds, independent of the
global number of runners.

\section*{Code and attribution}

The accompanying replay script verifies the two constructions and
the small affine rule exclusions. It checks local coverage and the
discrepancy in exact rational or integer arithmetic.
The source, checks and data are in the
\href{https://github.com/F1Justin/MathLab/tree/v1.0.0}{versioned MathLab repository}.
In the companion Lean project,
\texttt{MathLab.PeriodicPacking} verifies the counting period, sharp
discrepancy and the resulting integer implication. The geometric
construction uses the proofs and finite checker in this article.
Daria Poliakova devised the underlying three rules and the initial
framework \cite{Poliakova}. The paper's contributions are the specific
compatible compositions, the counting profiles and bounds, and the
finite reduction for the affine rule search. The minimum-frontier
search argument was developed with Astra.

Chengyi Fang initiated and funded the project, set its objectives,
supervised its execution and directed dissemination. Codex and
Astra performed the principal mathematical development, computation
and writing. The contributor understands aspects of the research
but has not independently checked any part of the proofs.

\medskip\noindent
This work is licensed under
\href{https://creativecommons.org/licenses/by/4.0/}{CC BY 4.0}.
The original accompanying code is released under the MIT License.

\begin{thebibliography}{9}
\bibitem{BCS}
M.~Blanco, F.~Criado and F.~Santos,
\emph{Coloopless zonotopes and counterexamples to the Shifted Lonely
Runner Conjecture}, \href{https://arxiv.org/abs/2603.24784v4}
{arXiv:2603.24784v4}, 2026.
\bibitem{Poliakova}
D.~Poliakova, \emph{More (shifted) runners, less loneliness},
\href{https://arxiv.org/abs/2609.23952v1}{arXiv:2609.23952v1}, 2026.
\bibitem{Zhang}
Y.~Zhang, \emph{Tight instances of the Lonely Runner Conjecture:
one-entry modifications and the growth of speeds},
\href{https://arxiv.org/abs/2608.13599v3}{arXiv:2608.13599v3}, 2026.
\bibitem{Upstream}
\texttt{aadrianvv} and \texttt{tomdif},
\emph{Shifted Lonely Runner counterexamples for \(18\le n\le57\)},
\href{https://github.com/1337isnot1337/solved-conjectures/issues/110}
{solved-conjectures, issue 110}, 2026.
\end{thebibliography}
\end{document}
